% ============================================================
%  Поглавље 5: QUBO — Quadratic Unconstrained Binary Optimization
% ============================================================
\section{QUBO: Quadratic Unconstrained Binary Optimization}
\label{sec:qubo}

Пре него што пређемо на квантне алгоритме, неопходно је увести
\emph{QUBO} формулацију. QUBO је стандардни математички оквир у коме
се комбинаторни оптимизациони проблеми изражавају на начин
директно погодан за квантне и квантно-инспирисане решаваче.

% --- Дефиниција ---
\subsection{Дефиниција и основна структура}

\begin{definition}[QUBO]
\emph{Проблем квадратне неограничене бинарне оптимизације}
(Quadratic Unconstrained Binary Optimization, QUBO) тражи
бинарни вектор $\mathbf{x} \in \{0,1\}^n$ који минимизује
(или максимизује) квадратну форму:
\[
  f(\mathbf{x})
  = \mathbf{x}^\top Q\, \mathbf{x}
  = \sum_{i=1}^{n} Q_{ii}\, x_i
    + 2\sum_{i < j} Q_{ij}\, x_i x_j,
\]
где је $Q \in \mathbb{R}^{n \times n}$ симетрична (или горња
троугаона) матрица коефицијената.
\end{definition}

\begin{example}[QUBO матрица за $n = 3$]
Нека је $n = 3$ и
\[
  Q =
  \begin{pmatrix}
    -1 &  1 &  0 \\
     1 & -1 &  1 \\
     0 &  1 & -1
  \end{pmatrix}.
\]
Тада је циљна функција
\[
  f(\mathbf{x}) = -x_1 - x_2 - x_3 + 2x_1 x_2 + 2x_2 x_3.
\]
За $\mathbf{x} = (1, 0, 1)^\top$ добијамо $f = -1 - 0 - 1 + 0 + 0 = -2$,
а за $\mathbf{x} = (1, 1, 1)^\top$ добијамо $f = -1 - 1 - 1 + 2 + 2 = 1$.
Минимум је дакле у $\mathbf{x}^* = (1, 0, 1)^\top$ са вредношћу $f^* = -2$.
\end{example}

Три кључне особине чине QUBO централним форматом:
\begin{itemize}[nosep]
  \item \textbf{Бинарност:} $x_i^2 = x_i$ за $x_i \in \{0,1\}$,
        па дијагонални чланови $Q_{ii} x_i^2$ редукују на
        линеарне чланове $Q_{ii} x_i$.
  \item \textbf{Без ограничења:} Сва ограничења се уграђују у
        циљну функцију преко казнених чланова (\emph{penalty terms}).
  \item \textbf{Квантна еквиваленција:} QUBO директно одговара
        \emph{Изинговом} (Ising) спин моделу, природном Хамилтонијану
        квантних анилера и QAOA алгоритма.
\end{itemize}

% --- Примери ограничења ---
\subsection{Примери QUBO формулација ограничења}

Конвертовање ограничених проблема у QUBO захтева да свако
ограничење заменимо са казненим чланом пропорционалним са
$\lambda > 0$ (довољно великим да казни свако кршење услова):

\textbf{Ограничење бинарности} $x_i + x_j = 1$ (тачно један
од два бита је 1):
\[
  P_1 = \lambda\,(x_i + x_j - 1)^2
      = \lambda\,(x_i + x_j - 2x_ix_j - 1 + 2x_ix_j)
      = \lambda\,(1 - x_i - x_j + 2x_ix_j - 1 + ...)
\]
Развијањем добијамо:
\[
  P_1 = \lambda\,\bigl(1 - x_i - x_j + 2x_ix_j\bigr).
\]
У QUBO матрици то одговара: $Q_{ii} \mathrel{-}= \lambda$,
$Q_{jj} \mathrel{-}= \lambda$, $Q_{ij} \mathrel{+}= 2\lambda$.

\textbf{Ограничење неједнакости} $x_i + x_j \leq 1$
(не могу оба бити 1 истовремено):
\[
  P_2 = \lambda\, x_i x_j.
\]
Кршење (оба $= 1$) додаје $\lambda$ циљној функцији.

\textbf{Тачно-$k$-од-$n$ ограничење} $\sum_{i=1}^n x_i = k$:
\[
  P_3 = \lambda\,\Bigl(\sum_{i=1}^n x_i - k\Bigr)^2
      = \lambda\,\Bigl(k^2 - 2k\sum_i x_i + \sum_i x_i^2 + 2\sum_{i<j}x_ix_j\Bigr).
\]
Пошто $x_i^2 = x_i$:
\[
  P_3 = \lambda\,\Bigl((1-2k)\sum_i x_i + 2\sum_{i<j}x_ix_j + k^2\Bigr).
\]

\begin{example}[3-бојење чворова — QUBO казна]
За ограничење да суседни чворови $u$ и $v$ у графу не смеју
добити исту боју $c$ уводимо бинарне варијабле $x_{u,c} \in \{0,1\}$.
Ограничење „$u$ и $v$ не деле боју $c$" кодира се казном:
\[
  P_{u,v,c} = \lambda\, x_{u,c}\, x_{v,c}.
\]
Свако кршење (оба чвора добијају боју $c$) повећава вредност
функције за $\lambda$, чиме решавач избегава та стања.
\end{example}

Табела~\ref{tab:qubo_constraints} сажима најчешће типове ограничења
и одговарајуће QUBO казнене чланове.

\begin{table}[ht]
\centering
\caption{Стандардна ограничења и њихове QUBO казне ($\lambda > 0$).}
\label{tab:qubo_constraints}
\renewcommand{\arraystretch}{1.6}
\resizebox{\linewidth}{!}{%
\begin{tabular}{|l|l|l|}
\hline
\textbf{Ограничење} & \textbf{Казнени члан $P$} & \textbf{Значење} \\
\hline
$x_i = 0$ & $\lambda\, x_i$ & Искључи чвор/ресурс из решења \\
\hline
$x_i = 1$ & $\lambda\,(1 - x_i)$ & Обавезно укључи у решење \\
\hline
$x_i + x_j = 1$ & $\lambda\,(1 - x_i - x_j + 2x_ix_j)$ & Тачно jedan од два бита активан \\
\hline
$x_i + x_j \leq 1$ & $\lambda\, x_i x_j$ & Не могу оба бити активна истовремено \\
\hline
$x_i = x_j$ & $\lambda\,(x_i + x_j - 2x_ix_j)$ & Оба бита морају имати исту вредност \\
\hline
$x_i \leq x_j$ & $\lambda\,(x_i - x_ix_j)$ & Ако је $x_i$ активан, мора бити и $x_j$ \\
\hline
$\sum_{i=1}^n x_i = k$ & $\lambda\,\left(\sum_{i=1}^n x_i - k\right)^2$ & Тачно $k$ битова активно \\
\hline
\end{tabular}}
\end{table}

\subsection{Пример примене QUBO казни: тачно-2-од-3 ограничење}

Нека су дате три бинарне променљиве $x_1, x_2, x_3 \in \{0,1\}$ и
линеарна циљна функција:
\[
  f_{\mathrm{obj}}(\mathbf{x}) = -4x_1 - 3x_2 - 5x_3.
\]
Желимо да минимизујемо $f_{\mathrm{obj}}$ уз ограничење да тачно две
од три промeнљиве морају бити активне:
\[
  x_1 + x_2 + x_3 = 2.
\]

\textbf{Казнени члан.}
Из табеле~\ref{tab:qubo_constraints}, ограничење $\sum_i x_i = k$
кодира се казном $\lambda\,(\sum_i x_i - k)^2$.
За $k = 2$ и $\lambda = 10$ добијамо:
\[
  P(\mathbf{x}) = 10\,(x_1 + x_2 + x_3 - 2)^2.
\]

\textbf{Потпуна QUBO формулација.}
\[
  f(\mathbf{x}) = f_{\mathrm{obj}}(\mathbf{x}) + P(\mathbf{x})
  = -4x_1 - 3x_2 - 5x_3 + 10\,(x_1 + x_2 + x_3 - 2)^2.
\]

\textbf{Провера на свим изводљивим решењима} (тачно два бита активна):
\begin{center}
\renewcommand{\arraystretch}{1.3}
\begin{tabular}{|c|c|c|c|}
\hline
$\mathbf{x}$ & $f_{\mathrm{obj}}$ & $P$ & $f$ \\
\hline
$(1,1,0)$ & $-7$ & $0$ & $-7$ \\
$(1,0,1)$ & $-9$ & $0$ & $\mathbf{-9}$ \\
$(0,1,1)$ & $-8$ & $0$ & $-8$ \\
\hline
\end{tabular}
\end{center}

За неизводљиво решење, нпр.\ $\mathbf{x} = (1,1,1)$:
\[
  f = -12 + 10\cdot(3-2)^2 = -12 + 10 = -2,
\]
што је лошије од сваког изводљивог решења, па решавач природно бира
$\mathbf{x}^* = (1, 0, 1)$ са вредношћу $f^* = -9$.

% --- MaxCut као QUBO ---
\subsection{Формулација MaxCut-а као QUBO}

Вратимо се проблему \textsc{MaxCut} из поглавља~\ref{sec:maxcut}.
Нека је $G = (V, E, w)$ граф са тежинама, $|V| = n$.

\textbf{Варијабле.}
Уведемо бинарну индикаторску променљиву за сваки чвор:
\[
  x_v \in \{0, 1\}, \quad v \in V,
\]
где $x_v = 0$ значи $v \in S$, а $x_v = 1$ значи $v \in \bar{S}$.

\textbf{Циљна функција.}
Ивица $\{u, v\}$ прелази рез ако и само ако $x_u \neq x_v$, тј.\
$x_u + x_v - 2x_u x_v = 1$.
Тежина реза коју \emph{максимизујемо} је:
\[
  \mathrm{cut}(\mathbf{x})
  = \sum_{\{u,v\} \in E} w_{uv}\,(x_u + x_v - 2x_u x_v).
\]

\textbf{Конверзија у QUBO минимизацију.}
Пошто QUBO стандардно тражи минимум, негирамо циљ:
\[
  \boxed{
  f_{\mathrm{MaxCut}}(\mathbf{x})
  = -\sum_{\{u,v\} \in E} w_{uv}\,(x_u + x_v - 2x_u x_v).
  }
\]

\textbf{Матрица $Q$.}
Елементи QUBO матрице читају се директно из горњег израза:
\[
  Q_{vv} = -\sum_{u:\{u,v\}\in E} w_{uv}
  \quad\text{(дијагонала)},
  \qquad
  Q_{uv} = w_{uv}
  \quad\text{(ванд‌ијагонала, $u \neq v$)}.
\]
Нема додатних ограничења. Проблем је \emph{природно неограничен},
па не треба казнени чланови.

\begin{example}[QUBO матрица за $K_4$ са јединичним тежинама]
За $K_4$ ($n = 4$, свих $6$ ивица са $w = 1$):
\[
  Q =
  \begin{pmatrix}
   -3 &  1 &  1 &  1 \\
    1 & -3 &  1 &  1 \\
    1 &  1 & -3 &  1 \\
    1 &  1 &  1 & -3
  \end{pmatrix}
\]
($Q_{vv} = -(n-1) = -3$ за сваки чвор,
$Q_{uv} = Q_{vu} = 1$ за сваку ивицу).

Оптималан бинарни вектор, нпр.\ $\mathbf{x} = (0,0,1,1)$
($S = \{0,1\}$, $\bar S = \{2,3\}$), даје:
\[
  f(\mathbf{x}) = -4 = -\mathrm{OPT}(K_4),
\]
у складу са резултатом из поглавља~\ref{sec:maxcut}.
\end{example}


\begin{remark}
QUBO је „заједнички језик" класичних и квантних оптимизационих
решавача: иста матрица $Q$ може бити прослеђен D-Wave анилеру,
класичном симулираном анилеру или QAOA алгоритму, без
икакве измене формулације проблема.
\end{remark}

% --- Решавање QUBO проблема ---
\subsection{Решавање QUBO: Z3 и симулирано жарење}

Да бисмо илустровали портабилност QUBO формата, решавамо исту
инстанцу MaxCut на $K_4$ са QUBO матрицом из претходног примера.
Решавамо потпуно различитим алатима: егзактним SMT решавачем Z3
и класичним симулираним жарењем преко D-Wave Ocean SDK-а.

\subsubsection{Решавање QUBO помоћу Z3}

Z3-ов \texttt{Optimize} пакет директно прима квадратну
циљну функцију над целобројним (или буловим) променљивама.
Бинарност се наноси као скуп ограничења $x_i \in \{0, 1\}$,
а решавач налазе егзактни минимум применом $\nu$Z оптимизационог оквира.

\begin{example}[QUBO за MaxCut на $K_4$ помоћу Z3 (Python)]
\begin{lstlisting}[style=python, caption={Решавање QUBO помоћу Z3 Optimize сучеља.}]
from z3 import *

# QUBO matrica za MaxCut na K4 (gornja-trougaona forma)
Q = {
    (0,0): -3, (1,1): -3, (2,2): -3, (3,3): -3,
    (0,1):  2, (0,2):  2, (0,3):  2,
    (1,2):  2, (1,3):  2,
    (2,3):  2,
}

n = 4
x = [Int(f'x{i}') for i in range(n)]
opt = Optimize()

# Binarne promenljive: x[i] in {0, 1}
for v in x:
    opt.add(Or(v == 0, v == 1))

# QUBO ciljna funkcija: sum_{i<=j} Q[i,j] * x[i] * x[j]
obj = Sum([
    Q.get((i, j), 0) * x[i] * x[j]
    for i in range(n) for j in range(i, n)
])
opt.minimize(obj)

if opt.check() == sat:
    m = opt.model()
    vals = [m[v].as_long() for v in x]
    print(f"Minimalna vrednost: {m.eval(obj)}")  # -4
    print(f"Resenje x: {vals}")                   # [0, 0, 1, 1]
\end{lstlisting}
\end{example}

Z3 проналази оптималан вектор $\mathbf{x} = (0,0,1,1)$ са вредношћу
$f(\mathbf{x}) = -4$, у складу са аналитичким резултатом.
Решавач гарантује глобалну оптималност, али временска сложеност
у најгорем случају остаје експоненцијална.

\textbf{Инверзни проблем: реконструкција матрице $Q$ из познатог решења.}
Сада разматрамо обрнут смер: претпоставимо да је познато да је
$\mathbf{x}^* = (0,0,1,1)$ \emph{оптимално} решење неког QUBO проблема,
али матрица $Q$ није позната.
Циљ је употребити Z3 да нађе $Q$ такву да $\mathbf{x}^*$ минимизује
$f(\mathbf{x}) = \mathbf{x}^\top Q\,\mathbf{x}$ над свим
$\mathbf{x} \in \{0,1\}^n$.

\textbf{Услов оптималности.}
Вектор $\mathbf{x}^*$ је глобални минимум ако и само ако:
\[
  f(\mathbf{x}^*) \leq f(\mathbf{x})
  \quad \forall\,\mathbf{x} \in \{0,1\}^n.
\]
За фиксно $\mathbf{x}^*$ и фиксно $\mathbf{x}$, ова неједнакост
је \emph{линеарна} у елементима матрице $Q$, стога је проблем
налажења $Q$ систем линеарних ограничења, природан за Z3.

\textbf{Неодређеност.}
Систем са $2^n$ ограничења и $n(n+1)/2$ слободних улаза у $Q$ је
генерално недоређен: много различитих матрица $Q$ може имати исти
оптимум. Додатна нормализациона ограничења (нпр.\ фиксирање скале
или структуре) отклањају ту вишезначност.

\begin{example}[Реконструкција $Q$ за коју је $(0,0,1,1)$ оптимум (Python)]
\begin{lstlisting}[style=python, caption={Z3 реконструкција QUBO матрице из познатог оптималног решења.}]
from z3 import *
from itertools import product

n      = 4
x_star = [0, 0, 1, 1]   # poznato optimalno resenje

# Slobodne (cjelobrojne) promenljive za gornje-trougaonu Q matricu
Q = {(i, j): Int(f'Q_{i}_{j}')
     for i in range(n) for j in range(i, n)}

def energy(x):
    """f(x) = sum_{i<=j} Q[i,j] * x[i] * x[j]"""
    return Sum([Q[(i, j)] * x[i] * x[j]
                for i in range(n) for j in range(i, n)])

s = Solver()

# Uslov optimalnosti: f(x*) <= f(x) za svaki x in {0,1}^n
f_star = energy(x_star)
for x in product([0, 1], repeat=n):
    s.add(f_star <= energy(list(x)))

# Normalizacija: fiksiramo offdijagonalne koeficijente na 2
# (odgovara jedinicnim tezinama MaxCut grafa)
for i in range(n):
    for j in range(i + 1, n):
        s.add(Q[(i, j)] == 2)

if s.check() == sat:
    m = s.model()
    print("Rekonstruisana Q matrica (gornje-trougaona):")
    for i in range(n):
        for j in range(i, n):
            print(f"  Q[{i},{j}] = {m[Q[(i, j)]]}")
\end{lstlisting}
\end{example}

З3 проналази јединствено решење под датим нормализационим
ограничењима:

\[
  Q_{ii} = -3 \;\text{(за сваки }i\text{)},
  \qquad
  Q_{ij} = 2 \;\text{(за сваки пар }i < j\text{)},
\]

чиме се добија тачно QUBO матрица $K_4$ MaxCut проблема.
Ово потврђује кохеренцију формулације: оптималност
$(0,0,1,1)$ у потпуности одређује структуру $Q$ (уз задату скалу).

\subsubsection{Решавање QUBO симулираним жарењем (D-Wave Ocean SDK)}

D-Wave Ocean SDK садржи класични \texttt{SimulatedAnnealingSampler}
који прима исти QUBO формат и враћа скуп узорака сортираних
по енергији. Самплер је квантно-инспирисан: не захтева стварни
квантни хардвер, али прати исту конвенцију интерфејса као и
правилни D-Wave анилер, тако да је прелаз на хардвер тривијалан
(замена самплера са \texttt{DWaveSampler}).

\begin{example}[QUBO за MaxCut на $K_4$ помоћу SimulatedAnnealingSampler (Python)]
\begin{lstlisting}[style=python, caption={Решавање QUBO симулираним жарењем преко D-Wave Ocean SDK-а.}]
from dwave.samplers import SimulatedAnnealingSampler

# Isti QUBO format kao i za Z3 (gornja-trougaona forma)
Q = {
    (0,0): -3, (1,1): -3, (2,2): -3, (3,3): -3,
    (0,1):  2, (0,2):  2, (0,3):  2,
    (1,2):  2, (1,3):  2,
    (2,3):  2,
}

sampler = SimulatedAnnealingSampler()
# num_reads: broj nezavisnih pokretanja zarenja
response = sampler.sample_qubo(Q, num_reads=100)

best = response.first
print(f"Energija: {best.energy}")       # -4.0
print(f"Resenje: {dict(best.sample)}")  # {0:0, 1:0, 2:1, 3:1}
\end{lstlisting}
\end{example}

Метод \texttt{sample\_qubo} покреће \texttt{num\_reads} независних
трагова жарења и враћа резултате у облику \texttt{SampleSet} објекта.
Атрибут \texttt{response.first} садржи узорак са најнижом енергијом.
За прелаз на стварни D-Wave хардвер довољно је заменити самплер:

\begin{lstlisting}[style=python, caption={Прелаз на D-Wave хардвер тј. само се самплер мења.}]
from dwave.system import DWaveSampler, EmbeddingComposite

sampler = EmbeddingComposite(DWaveSampler())
response = sampler.sample_qubo(Q, num_reads=100)
\end{lstlisting}

\texttt{EmbeddingComposite} аутоматски проналази уграђивање
логичких QUBO чворова у физички граф процесора,
скривајући детаље архитектуре од корисника.


\begin{remark}
Оба приступа прихватају идентичан QUBO формат $Q$.
Z3 гарантује егзактно решење, али не скалира на велике инстанце.
Симулирано жарење је хеуристичко и не гарантује оптималност,
али лако обрађује инстанце са стотинама варијабли.
Квантни анилер (D-Wave) нуди исти интерфејс уз потенцијалну
предност тунелирања на хардверском нивоу.
\end{remark}
